%%% Mandatory information page of the thesis


\null
\vfill
\vfill
\vfill

\begin{center}
    {\Large\bfseries Abstract}
\end{center}

\vfill

Positivity bounds provide a powerful framework for constraining effective field theories using general principles such as unitarity, locality, crossing symmetry and causality. In this work, we investigate how different assumptions about the ultraviolet behaviour of scattering amplitudes restrict the low-energy spectrum and the structure of effective field theories.

We first study positivity bounds in theories of massive and massless scalars satisfying different subtraction schemes in dispersion relations, deriving analytic and numerical constraints on the allowed spectrum and resonance couplings. In particular, we show that theories satisfying $0$-subtracted dispersion relations require the presence of a lighter scalar resonance in order to consistently accommodate a heavier spin-$2$ state.

We then develop a systematic framework for constructing scattering amplitudes of non-Abelian massive vectors and deriving positivity bounds both in the forward limit and at finite momentum transfer.
Finally, we apply these methods to electroweak $WW$ scattering and effective theories beyond the Standard Model. We analyze how positivity bounds constrain deformations of the electroweak sector and the contributions of higher-order operators under different subtraction schemes.

\vfill
\vfill
\vfill
\null